\appendices
\section{Reproducibility}\label{sec:appendix}

\subsection{Architecture Diagram}
The full-width architecture diagram is rendered in
Figure~\ref{fig:arch} (Section~\ref{sec:arch}). The dataflow
companion (\texttt{paper/figs/dataflow.eps}) traces a single
transaction from FFI entry through the WAL fsync; we include it in
the paper artifact rather than reprinting it here.

\subsection{Bench Command Lines}
The certification run that produced every number in this paper was
launched via
\begin{lstlisting}[language={}]
xbabe1_run.sh certify \
  --engines redline,sqlite \
  --workloads point-read-pk,point-read-secondary,\
secondary-index-range,mixed95-read5-write,\
mixed50-read50-write,writers-disjoint,\
hot-row-update,insert-only \
  --durabilities strict,relaxed \
  --threads 1,2,4,8,16,32,64,128 \
  --reps 5 --warmup 1 \
  --reserve-cores 4 \
  --emit target/bench/xbabe1/certification
\end{lstlisting}
Recovery and failpoint matrices are launched separately:
\begin{lstlisting}[language={}]
xbabe1_run.sh recovery --cases all
xbabe1_run.sh failpoint --cases all
xbabe1_run.sh compat --cases all
\end{lstlisting}

\subsection{Manifest Excerpt}
Each certification run emits a \texttt{manifest.json} pinning the
git SHA, the Docker image digest, the host fingerprint, the SQLite
PRAGMA snapshot, and per-output checksums:
\begin{lstlisting}[language={}]
{
  "git_sha": "<commit>",
  "docker_image_digest": "sha256:<digest>",
  "host": {
    "cpu": "x86_64 128 vCPU",
    "kernel": "Linux 6.8",
    "fs": "ext4",
    "docker": "29.2.1",
    "rustc": "1.95.0",
    "rusqlite": "0.37"
  },
  "sqlite_pragmas": {
    "journal_mode": "WAL",
    "synchronous_strict": "FULL",
    "synchronous_relaxed": "NORMAL"
  },
  "checksums": { "...": "..." }
}
\end{lstlisting}
The full manifest is reproduced in the artifact under the
certification output directory referenced above.
